\section*{Further Reading}
\addcontentsline{toc}{section}{Further Reading}
\label{sec:further-reading}

\begin{layman}
This section is for readers who want to dig further into the
mathematics or formalization tools that this document touches.
Nothing here is cited in the body --- these are simply pointers to
standard textbooks, open-access notes, and lecture series that cover
the same territory at greater length and detail. References are
grouped by topic; difficulty ranges from advanced undergraduate to
research-level, and we have flagged the more elementary entries.
\end{layman}

\subsection*{Riemann surfaces}

\begin{itemize}
  \item Otto Forster, \emph{Lectures on Riemann Surfaces}. Springer
    Graduate Texts in Mathematics 81, 1981. The standard concise
    graduate textbook on the analytic side.
  \item Rick Miranda, \emph{Algebraic Curves and Riemann Surfaces}.
    AMS Graduate Studies in Mathematics 5, 1995. The most
    accessible introduction; bridges the analytic and algebraic
    viewpoints throughout.
  \item Hershel M.\ Farkas and Irwin Kra, \emph{Riemann Surfaces}.
    Springer Graduate Texts in Mathematics 71, 2nd ed.\ 1992.
    Encyclopedic; covers theta functions, Jacobians, automorphism
    groups in depth.
  \item Simon K.\ Donaldson, \emph{Riemann Surfaces}. Oxford
    Graduate Texts in Mathematics 22, 2011. Modern, geometric, and
    well-illustrated.
\end{itemize}

\subsection*{Algebraic curves and algebraic geometry}

\begin{itemize}
  \item Phillip Griffiths and Joseph Harris, \emph{Principles of
    Algebraic Geometry}. Wiley, 1978. The classic reference for
    the Hodge-theoretic and complex-analytic perspective.
  \item Robin Hartshorne, \emph{Algebraic Geometry}. Springer
    Graduate Texts in Mathematics 52, 1977. The standard
    scheme-theoretic textbook.
  \item Enrico Arbarello, Maurizio Cornalba, Phillip Griffiths, and
    Joseph Harris, \emph{Geometry of Algebraic Curves}, Volume I.
    Springer Grundlehren 267, 1985. Definitive reference for moduli
    of curves and the Brill--Noether theory underpinning the
    Abel--Jacobi map.
  \item David Mumford, \emph{Curves and Their Jacobians}. University
    of Michigan Press, 1975. Short, lucid lecture notes.
\end{itemize}

\subsection*{Jacobians, abelian varieties, and the Abel--Jacobi map}

\begin{itemize}
  \item Christina Birkenhake and Herbert Lange, \emph{Complex Abelian
    Varieties}. Springer Grundlehren 302, 2nd ed.\ 2004. The
    standard reference for the analytic theory.
  \item David Mumford, \emph{Abelian Varieties}. Tata Institute /
    Oxford University Press, 1970. Algebraic-geometric companion to
    the analytic story.
  \item Alexander Polishchuk, \emph{Abelian Varieties, Theta
    Functions and the Fourier Transform}. Cambridge Tracts in
    Mathematics 153, 2003.
\end{itemize}

\subsection*{Complex analysis (one variable)}

\begin{itemize}
  \item Lars V.\ Ahlfors, \emph{Complex Analysis}. McGraw-Hill,
    3rd ed.\ 1979. The classical gold standard for one-variable
    complex analysis.
  \item Elias M.\ Stein and Rami Shakarchi, \emph{Complex Analysis}.
    Princeton Lectures in Analysis II, Princeton University Press,
    2003. Modern and very readable.
\end{itemize}

\subsection*{Smooth manifolds, differential geometry, Lie groups}

\begin{itemize}
  \item John M.\ Lee, \emph{Introduction to Smooth Manifolds}.
    Springer Graduate Texts in Mathematics 218, 2nd ed.\ 2012.
    The standard graduate textbook on smooth manifolds.
  \item Frank W.\ Warner, \emph{Foundations of Differentiable
    Manifolds and Lie Groups}. Springer Graduate Texts in
    Mathematics 94, 1983. Includes a clean treatment of de Rham
    cohomology and the Hodge theorem on a compact Riemannian
    manifold.
\end{itemize}

\subsection*{Algebraic topology and homology}

\begin{itemize}
  \item Allen Hatcher, \emph{Algebraic Topology}. Cambridge
    University Press, 2002. Free PDF from the author's homepage:
    \url{https://pi.math.cornell.edu/~hatcher/AT/ATpage.html}.
  \item Raoul Bott and Loring W.\ Tu, \emph{Differential Forms in
    Algebraic Topology}. Springer Graduate Texts in Mathematics 82,
    1982. The cleanest entry into de Rham theory and spectral
    sequences for working geometers.
  \item William S.\ Massey, \emph{A Basic Course in Algebraic
    Topology}. Springer Graduate Texts in Mathematics 127, 1991.
    Conventional, classical exposition; useful as a second source.
\end{itemize}

\subsection*{Hodge theory, Sobolev spaces, elliptic regularity}

\begin{itemize}
  \item Raymond O.\ Wells, Jr., \emph{Differential Analysis on
    Complex Manifolds}. Springer Graduate Texts in Mathematics 65,
    3rd ed.\ 2008. The standard book-length treatment of Hodge
    theory and Dolbeault cohomology on K\"ahler manifolds.
  \item Jean-Pierre Demailly, \emph{Complex Analytic and
    Differential Geometry}. Open-access manuscript, available from
    the author's webpage at Universit\'e Grenoble Alpes:
    \url{https://www-fourier.ujf-grenoble.fr/~demailly/manuscripts/agbook.pdf}.
    Comprehensive reference for sheaf theory, $L^2$ methods, and
    Hodge theory.
  \item Lawrence C.\ Evans, \emph{Partial Differential Equations}.
    AMS Graduate Studies in Mathematics 19, 2nd ed.\ 2010. The
    standard graduate PDE textbook; covers Sobolev spaces and
    elliptic regularity at the level needed for Hodge theory.
  \item Michael E.\ Taylor, \emph{Partial Differential Equations
    I--III}. Springer Applied Mathematical Sciences 115--117. The
    pseudodifferential and spectral chapters are the natural
    classical companions to the R10 framework discussed in
    Section~12.
\end{itemize}

\subsection*{Lean 4, Mathlib, and formalized mathematics}

\begin{itemize}
  \item Jeremy Avigad, Leonardo de Moura, Soonho Kong, and Sebastian
    Ullrich, \emph{Theorem Proving in Lean 4}. Online textbook:
    \url{https://leanprover.github.io/theorem_proving_in_lean4/}.
  \item Jeremy Avigad and Patrick Massot, \emph{Mathematics in
    Lean}. Online tutorial:
    \url{https://leanprover-community.github.io/mathematics_in_lean/}.
    Walks through formalizing real mathematics in current Mathlib.
  \item Mathlib documentation (the library this project is built
    against): \url{https://leanprover-community.github.io/mathlib4_docs/}.
  \item Kevin Buzzard, \emph{The Xena Project} blog:
    \url{https://xenaproject.wordpress.com/}. Long-running record of
    formalized-mathematics work and pedagogy from one of the
    field's most active organizers --- and the author of the
    challenge this project addresses.
\end{itemize}

\subsection*{Talks and lecture videos}

\begin{itemize}
  \item Richard E.\ Borcherds runs an extensive YouTube lecture
    series covering graduate algebraic geometry, schemes, number
    theory, and Lie groups, accessible to mathematics graduate
    students.
  \item Kevin Buzzard has given many public talks on formalized
    mathematics, the Xena Project, and the formalization of
    Fermat's Last Theorem; recordings are easy to find on YouTube.
  \item 3Blue1Brown (Grant Sanderson):
    \url{https://www.youtube.com/@3blue1brown}. Visual introductions
    to topology, complex analysis, and differential equations,
    aimed at a general mathematically-curious audience --- a good
    on-ramp for the plain-English readers of this document.
\end{itemize}
